Automatic sleep staging models can lose accuracy when applied to unseen
patients, yet a single label-free remedy is usually applied: replace the
normalization statistics of the target recording. This assumes that subject
shift is a recording-wide offset. We introduce a patient-identity diagnostic
that tests this assumption by refitting a linear subject probe within each
sleep stage. After marginal mean correction, the within-stage probe remains
well above chance, whereas the overall probe is dominated by the removable
offset. We therefore propose stage-conditional shift normalization (SCSN), a
closed-form, label-free correction that estimates a separate shift for each
stage from the model's soft posterior and subtracts the posterior-weighted
shift during inference. SCSN contains marginal correction as its $\lambda=0$
endpoint and applies unchanged to batch- and group-normalized networks.
Across Sleep-EDF (39 recordings), HMC (151 recordings), two backbones and two
normalization schemes, SCSN reduced the within-stage subject probe from
$0.74$ to $0.47$ and from $0.53$ to $0.21$ on EEGNet. Classification accuracy
was preserved on Sleep-EDF and improved by $1.4$--$1.9$ points on HMC. On
group-normalized models, SCSN significantly improved accuracy in three of four
settings and improved worst-patient conformal coverage in all four. Cross-cohort
transfer also exposes the method's precondition: a weak target posterior
reduces SCSN to marginal correction, whereas oracle stages recover most of the
gap.
